%% =========================================================================
%% Demolition-indicator theory section (2026-07-11 restructure).
%% Consolidates the former BBR subsection "Demolition indicators used in
%% prior register-based studies" and the former Methods subsection
%% "Candidate demolition indicators" into one section: it first introduces
%% the D1–D6 indicator set (from bbr-demolitions/src/indicators.py) with
%% schematics and definitions, then maps prior published studies onto it.
%% The indicator set is presented one indicator per subsubsection, in
%% numeric order D1–D6, each with its own schematic figure.
%% =========================================================================
\section{Demolition indicators}\label{sec:indicators}

No field in BBR directly flags a completed physical demolition
(Section~\ref{sec:bbr-chain}). Every demolition indicator -- whether used
in a prior published study or computed here -- is therefore a proxy: an
observable register event, sometimes combined with filters intended to
remove false positives.

Because the completion fields are not distributed
(Section~\ref{sec:bbr-chain}), this study defines a demolition
operationally from the two traces of the workflow that \emph{are}
observable: a building's own lifecycle-status and business-process
history (\texttt{Bygning}), and its linked demolition case
(\texttt{Sagsniveau}, \texttt{BBRSag}). This gives two families of
indicators, contiguously numbered: D1--D3 from the building's own
history, D4--D6 from its linked case. Section~\ref{sec:ind-set}
illustrates the proxy mechanisms and defines the resulting six indicators,
and Section~\ref{sec:lit-indicators} maps the set onto the choices made
in prior published Danish studies.

\subsection{Proxy mechanisms and the indicator set (D1--D6)}\label{sec:ind-set}

The six indicators rest on a small number of proxy mechanisms, each an
observable register event standing in for the unobserved physical demolition.
All six are computed on the same dataset, clipped to the full reporting
window (2000--2025 inclusive; Section~\ref{sec:data}) and are otherwise
unfiltered; a building that matches an indicator but lacks a date under
it (e.g., a linked case with no status-10 year) is kept rather than
dropped, so that intersection indicators are not silently altered. Each
indicator is presented in turn, D1 through D6, with its own schematic;
the intersection indicators (D3 and D5) combine primitive signals
defined just above them.

\subsubsection*{Register-exit proxy (Indicator D1)}

The register-exit signal reads the building's own history to identify
when it leaves the active stock (Figure~\ref{fig:ind-d1}). D1 counts a
building as demolished if its status ever became 10 (``Historisk''); the
year is the earliest such effect date. This is the maximally inclusive
register-exit signal: status 10 is also reached by re-registration,
merges, and error correction, so D1 overcounts by design, and the
current definition does not check for later reactivation to an active
status. Discontinued use codes are treated as a separate, composable
axis (Section~\ref{sec:discontinued-method}), not folded into D1 itself.

\begin{figure}[h!]
\centering
\begin{tikzpicture}[node distance=6mm and 8mm]
\node[ind] (a) {Building record is updated};
\node[ind, below=of a] (b) {Lifecycle status becomes 10, ``Historisk'' (or \texttt{ObjStatus}\,=\,``udg{\aa}et'')};
\node[ind, below=of b] (c) {Building treated as demolished};
\draw[indarrow] (a) -- (b);
\draw[indarrow] (b) -- (c);
\end{tikzpicture}
\caption{Register-exit proxy (indicator~D1): a building leaving the
active stock is treated as demolished. This is the signal behind the
register-exit studies in Table~\ref{tab:lit-indicators}.}
\label{fig:ind-d1}
\end{figure}

\subsubsection*{Event-coded update proxy (Indicator D2)}

This proxy relies entirely on the business-process code attached to a
building's update history (Figure~\ref{fig:ind-d2}). D2 counts a building
whose history ever carried \texttt{ForretningsProcess} 3 (``updated due
to demolition''); the year is the earliest such effect date. D2 is not a
subset of D1: on the current extract, 126{,}362 buildings carry process
3 without ever reaching status 10, so this code is missed by many real
demolitions that are registered through another workflow (undercount).

\begin{figure}[h!]
\centering
\begin{tikzpicture}[node distance=6mm and 8mm]
\node[ind] (a) {Building record is updated};
\node[ind, below=of a] (b) {Update is coded \texttt{ForretningsProcess}\,=\,3, ``updated due to demolition''};
\node[ind, below=of b] (c) {Building treated as demolished};
\draw[indarrow] (a) -- (b);
\draw[indarrow] (b) -- (c);
\end{tikzpicture}
\caption{Event-coded update proxy (indicator~D2): treats the
demolition-coded record update as the demolition signal.}
\label{fig:ind-d2}
\end{figure}

\subsubsection*{Register exit and demolition process (Indicator D3)}

This proxy fires only where the register exit (D1) coincides with the
demolition-coded update (D2) on the same building
(Figure~\ref{fig:ind-d3}). D3 is the intersection D1 $\cap$ D2, dated by
D1's register-exit year. Because D2 is not a subset of D1, D3 is markedly
smaller than either indicator alone.

\begin{figure}[h!]
\centering
\begin{tikzpicture}[node distance=16mm and 10mm]
\node[ind, text width=4.7cm] (a) {Register exit: status 10, ``Historisk'' (D1)};
\node[ind, text width=4.7cm, right=of a] (b) {Demolition-coded update: \texttt{ForretningsProcess}\,=\,3 (D2)};
\node[ind, text width=6.2cm, below=of $(a.south)!0.5!(b.south)$] (c) {Both signals present: building treated as demolished, dated by the register exit};
\draw[indarrow] (a.south) -- (c.north);
\draw[indarrow] (b.south) -- (c.north);
\end{tikzpicture}
\caption{Intersection proxy (indicator~D3): requires the register exit
(D1) to coincide with the demolition process code (D2).}
\label{fig:ind-d3}
\end{figure}

\subsubsection*{Total demolition case (Indicator D4)}

From here the proxies read the building's linked demolition case rather
than its status history. A case is tied to a building through the
\texttt{stamdataBygning} field on \texttt{Sagsniveau}, which names the
building the case is filed on; D4 keeps only total-demolition cases
(Figure~\ref{fig:ind-d4}). D4 counts a building linked, via
\texttt{Sagsniveau}, to a demolition case of type
\texttt{sagstype}\,=\,32 (``Nedrivning hel,'' total). Partial-demolition
cases (\texttt{sagstype}\,=\,31, ``Nedrivning delvis'') also exist in the
register but qualify no indicator in this study, since a partially
demolished building generally remains standing. The year is the
building's own status-10 date where available, and left undated
otherwise. Roughly 82{,}000 demolition-case rows carry no linkable
building (\texttt{stamdataBygning} null) and are necessarily excluded, an
undercount shared with D6.

\begin{figure}[h!]
\centering
\begin{tikzpicture}[node distance=6mm and 8mm]
\node[ind] (a) {Demolition case linked to the building (\texttt{Sagsniveau}, \texttt{stamdataBygning})};
\node[ind, below=of a] (b) {Case type is \texttt{sagstype}\,=\,32, ``Nedrivning hel'' (total)};
\node[ind, below=of b] (c) {Building treated as demolished; dated by its status-10 exit where available};
\draw[indarrow] (a) -- (b);
\draw[indarrow] (b) -- (c);
\end{tikzpicture}
\caption{Demolition-case proxy, total (indicator~D4): a building linked
to a total-demolition case is treated as demolished, dated by its
register exit.}
\label{fig:ind-d4}
\end{figure}

\subsubsection*{Register exit and total demolition case (Indicator D5)}

This proxy requires the register exit (D1) to coincide with a linked
total-demolition case (D4) on the same building
(Figure~\ref{fig:ind-d5}). D5 is the intersection D1 $\cap$ D4, dated by
D1's register-exit year: a building that both left the active stock and
has a linked total-demolition case.

\begin{figure}[h!]
\centering
\begin{tikzpicture}[node distance=16mm and 10mm]
\node[ind, text width=4.7cm] (a) {Register exit: status 10, ``Historisk'' (D1)};
\node[ind, text width=4.7cm, right=of a] (b) {Linked total-demolition case: \texttt{sagstype}\,=\,32 (D4)};
\node[ind, text width=6.2cm, below=of $(a.south)!0.5!(b.south)$] (c) {Both signals present: building treated as demolished, dated by the register exit};
\draw[indarrow] (a.south) -- (c.north);
\draw[indarrow] (b.south) -- (c.north);
\end{tikzpicture}
\caption{Intersection proxy (indicator~D5): requires the register exit
(D1) to coincide with a linked total-demolition case (D4).}
\label{fig:ind-d5}
\end{figure}

\subsubsection*{Case-date proxy for felt 295 (Indicator D6)}

D6 dates a linked total-demolition case by its notification date,
standing in for the undistributed completion field felt 295
(Figure~\ref{fig:ind-d6}). D6 counts a building linked to a
total-demolition case (\texttt{sagstype}\,=\,32, as in D4) with a non-null
case notification date (\texttt{sag002 Byggesagsdato}, the approximate
counterpart of the undistributed felt 294 notification field); the year
is that date. Felt 295 itself is not reconstructable from this extract
(Section~\ref{sec:bbr-chain}): the one candidate completion field carried
by the case data (\texttt{sag010 FuldførelseAfByggeri}) is both sparse
and typically dates a co-filed construction case rather than the
demolition, so it is not used. D6 therefore dates demolition
\emph{notification}, not confirmed completion.

\begin{figure}[h!]
\centering
\begin{tikzpicture}[node distance=6mm and 8mm]
\node[ind] (a) {Total-demolition case linked to the building (\texttt{sagstype}\,=\,32)};
\node[ind, below=of a] (b) {Case carries a notification date (\texttt{sag002 Byggesagsdato}, $\approx$ felt 294)};
\node[ind, below=of b] (c) {Building treated as demolished; dated by the notification date};
\draw[indarrow] (a) -- (b);
\draw[indarrow] (b) -- (c);
\end{tikzpicture}
\caption{Case-date proxy (indicator~D6): a linked total-demolition case
is dated by its notification date, an approximation of the undistributed
felt~295 completion signal.}
\label{fig:ind-d6}
\end{figure}

None of D1--D6 is treated as a reference indicator in this study, and
none is scored against another as ground truth; Sections~\ref{sec:comparisons}
and~\ref{sec:results} describe how the six are compared and validated
instead.

\begin{table}[htbp]
\centering
\small
\caption{The demolition-indicator set computed in this study, unfiltered
building counts windowed to 2000--2025 (denominator: 6{,}250{,}075
buildings in the extract).}
\label{tab:d-indicators}
\begin{tabular}{@{}lp{7.4cm}r@{}}
\toprule
ID & Signal & Buildings \\
\midrule
D1 & status = 10 & 436{,}194 \\
D2 & forretningsproces = 3 & 225{,}706 \\
D3 & D1 $\cap$ D2 & 99{,}664 \\
D4 & sagstype = 32 (total) & 237{,}012 \\
D5 & D1 $\cap$ D4 & 200{,}634 \\
D6 & sagstype = 32 with notification date & 231{,}962 \\
\bottomrule
\end{tabular}
\end{table}

\subsection{Demolition indicators in prior register-based studies}\label{sec:lit-indicators}

Table~\ref{tab:lit-indicators} lists the approach and key filters used
by published Danish studies that estimate building demolition from BBR or
a BBR-derived extract, and, in its second column, the indicator in the
set above (Section~\ref{sec:ind-set}) that rests on the same register
signal. The register-exit and event-coded studies map directly onto D1
and D2. Two further approaches used in the literature have no direct
counterpart in the D1--D6 set, because both reach outside the register
signals this study computes: studies working from a pre-made, opaque
demolition extract, whose filtering recipe is not
published;\footnote{The recipe is unpublished, but the underlying signal is
recoverable. For the one such extract we hold --- the KMD/Andersen
extract~\cite{andersen2022adaptation} --- every case is a total demolition
(\texttt{sagstype = 32}), and the case-based indicator reproduces its
membership to $\sim$99\,\% window-matched (Section~\ref{sec:kmd}). So it
corresponds empirically to the case family D4--D6, even though the authors
state no register field that identifies it.} and change-detection studies,
which infer demolition from a teardown-and-rebuild sequence rather than
from a demolition signal on the building itself.

Two entries in Table~\ref{tab:lit-indicators} fall outside the proxy
scheme entirely. SBi's evaluation of village-renewal demolition
subsidies~\cite{jensen2019landsbyfornyelse} additionally draws on
BOSSINF, the administrative register of grant-funded demolitions: this
is genuine ground truth rather than a proxy, but limited to subsidised
cases in 54 eligible municipalities. Aagaard et
al.~\cite{aagaard2013levetider} do not detect individual demolitions at
all; they assume an aggregate demolition rate ($\approx$0.3\,\%/yr)
derived from construction-cost statistics, which underlies the standard
Danish service-life tables and is the baseline that later extract-based
studies benchmark against.

\begin{table}[h!]
\centering
\small
\renewcommand{\arraystretch}{1.2}
\setlength{\tabcolsep}{5pt}
\newcolumntype{L}[1]{>{\raggedright\arraybackslash}p{#1}}
\caption{Demolition indicators (proxies) and filters used in prior
published Danish register-based studies, and the corresponding indicator
in this study's set (Section~\ref{sec:ind-set}).}
\label{tab:lit-indicators}
\begin{tabular}{@{}L{0.25\linewidth}L{0.31\linewidth}L{0.34\linewidth}@{}}
\toprule
Study & Approach (corresponding indicator) & Key filters \\
\midrule
Aagaard et al.\ (2013)~\cite{aagaard2013levetider} &
  Not a detection method (n/a) &
  n/a \\
\addlinespace[2pt]\cmidrule(l{0pt}r{0pt}){1-3}\addlinespace[2pt]
{\O}stergaard et al.\ (2018)~\cite{ostergaard2018temporal} &
  Opaque extract (no direct D-indicator) &
  Missing/duplicate records dropped; lifespan $<5$\,yr dropped; built
  before 1800 dropped \\
\addlinespace[2pt]\cmidrule(l{0pt}r{0pt}){1-3}\addlinespace[2pt]
SBi (2019:01)~\cite{jensen2019landsbyfornyelse} &
  Register-exit (D1); plus BOSSINF ground truth &
  Residential codes only; 54 eligible municipalities \\
\addlinespace[2pt]\cmidrule(l{0pt}r{0pt}){1-3}\addlinespace[2pt]
Andersen \& Negendahl (2022)~\cite{andersen2022adaptation} &
  Opaque extract (no direct D-indicator) &
  Pre-2011 records discarded \\
\addlinespace[2pt]\cmidrule(l{0pt}r{0pt}){1-3}\addlinespace[2pt]
BUILD (2022:36)~\cite{jensen2022enfamiliehuse} &
  Change detection, BBR flag (outside D1--D6) &
  Code 120 only; same matrikel; rebuild after 2010 \\
\addlinespace[2pt]\cmidrule(l{0pt}r{0pt}){1-3}\addlinespace[2pt]
Andersen \& Negendahl (2023)~\cite{andersen2023lifespan} &
  Opaque extract (no direct D-indicator) &
  Missing date/area/use dropped; pre-2010 dropped; construction year
  ``1000'' (auto-generated) dropped \\
\addlinespace[2pt]\cmidrule(l{0pt}r{0pt}){1-3}\addlinespace[2pt]
Boligøkonomisk Videncenter (2023)~\cite{blv2023karakteristika} &
  Change detection, transaction (outside D1--D6) &
  Arms-length sales only; plot 500--2{,}000\,m$^2$; price
  15--15{,}000\,kr/m$^2$; matrikel age $\geq$20\,yr \\
\addlinespace[2pt]\cmidrule(l{0pt}r{0pt}){1-3}\addlinespace[2pt]
BUILD (2025)~\cite{jensen2025nedrivning} &
  Register-exit (D1) &
  Buildings built before 1999 only \\
\addlinespace[2pt]\cmidrule(l{0pt}r{0pt}){1-3}\addlinespace[2pt]
Droob et al.\ (2026)~\cite{droob2026servicelife} &
  Event-coded update (D2) &
  Outbuildings removed \\
\bottomrule
\end{tabular}
\end{table}
